\nivelbasico


\begin{ejercicio}
Un automóvil se encuentra en una carretera recta. Se elige un árbol como origen del sistema de referencia y se toma como sentido positivo el que apunta hacia el este.

En cierto instante, el automóvil se encuentra a \(\SI{12}{\meter}\) al este del árbol.

\begin{enumerate}[label=\alph*)]
    \item ¿Cuál es la posición del automóvil?
    \item ¿Qué información proporciona el signo de la posición?
    \item ¿La posición indica por sí sola hacia dónde se mueve el automóvil?
\end{enumerate}

\textit{Pista: la posición indica dónde se encuentra el objeto respecto del origen; no describe necesariamente su movimiento.}
\end{ejercicio}

\begin{respuesta}
\begin{enumerate}[label=\alph*)]
    \item \(x=+\SI{12}{\meter}\).
    \item El signo positivo indica que el automóvil está en el lado positivo del origen.
    \item No.
\end{enumerate}
\end{respuesta}

\begin{explicacion}
El origen se encuentra en el árbol:

\[ x=0. \]

El sentido positivo fue definido hacia el este. Como el automóvil se encuentra \(\SI{12}{\meter}\) al este del árbol:

\[ \boxed{x=+\SI{12}{\meter}}. \]

El signo positivo no significa que el automóvil se esté moviendo hacia el este. Indica únicamente que su \textbf{posición} se encuentra en el lado positivo del sistema de referencia.

El automóvil podría estar:

\begin{itemize}
    \item moviéndose hacia el este;
    \item moviéndose hacia el oeste;
    \item momentáneamente en reposo.
\end{itemize}

Por tanto, conocer la posición de un objeto no basta para determinar su estado de movimiento.
\end{explicacion}


\begin{ejercicio}
Sobre una carretera recta se establece un origen \(O\) y se toma hacia la derecha como sentido positivo.

Tres objetos se encuentran en las posiciones mostradas:

\begin{center}
\begin{tikzpicture}[scale=0.9]

    \draw[-{Latex},thick] (-5.5,0) -- (6,0) node[right] {\(x\)};

    \foreach \x/\lab in {-4/-4,-2/-2,0/0,2/2,4/4}{
        \draw (\x,0.12) -- (\x,-0.12);
        \node[below] at (\x,-0.12) {\(\lab\)};
    }

    \fill (-4,0) circle (2.5pt);
    \node[above] at (-4,0.12) {\(A\)};

    \fill (0,0) circle (2.5pt);
    \node[above] at (0,0.12) {\(B\)};

    \fill (4,0) circle (2.5pt);
    \node[above] at (4,0.12) {\(C\)};

    \node[below] at (5,-0.35) {\(\si{\meter}\)};

\end{tikzpicture}
\end{center}

Determine la posición de cada objeto e interprete físicamente su signo.

\textit{Pista: el origen tiene posición cero; los puntos a su izquierda tienen coordenada negativa.}
\end{ejercicio}

\begin{respuesta}
\[ x_A=-\SI{4}{\meter},\qquad x_B=\SI{0}{\meter},\qquad x_C=+\SI{4}{\meter}. \]
\end{respuesta}

\begin{explicacion}
El origen del sistema corresponde a:

\[ x=0. \]

El punto \(A\) está a \(\SI{4}{\meter}\) hacia el lado negativo del eje:

\[ \boxed{x_A=-\SI{4}{\meter}}. \]

El objeto \(B\) se encuentra exactamente en el origen:

\[ \boxed{x_B=\SI{0}{\meter}}. \]

El objeto \(C\) está a \(\SI{4}{\meter}\) hacia el lado positivo:

\[ \boxed{x_C=+\SI{4}{\meter}}. \]

Los signos indican en qué lado del origen se encuentran los objetos. No representan una distancia negativa ni permiten determinar por sí solos hacia dónde se están moviendo.
\end{explicacion}


\begin{ejercicio}
Una partícula se encuentra inicialmente en:

\[ x_i=-\SI{4}{\meter}. \]

Después se encuentra en:

\[ x_f=+\SI{9}{\meter}. \]

Calcule su desplazamiento.

\textit{Pista: utilice \(\Delta x=x_f-x_i\).}
\end{ejercicio}

\begin{respuesta}
\[ \boxed{\Delta x=+\SI{13}{\meter}}. \]
\end{respuesta}

\begin{explicacion}
El desplazamiento se define como el cambio de posición:

\[ \Delta x=x_f-x_i. \]

Sustituyendo:

\[
\Delta x
=
\SI{9}{\meter}
-
(-\SI{4}{\meter}).
\]

Entonces:

\[ \boxed{\Delta x=+\SI{13}{\meter}}. \]

El signo positivo indica que el desplazamiento neto tiene el sentido positivo del eje.

Observe que el desplazamiento depende únicamente de la posición inicial y la posición final. No necesitamos conocer todavía cómo se movió la partícula entre ambos puntos.
\end{explicacion}


\begin{ejercicio}
Una persona camina en línea recta desde:

\[ x_i=\SI{2}{\meter} \]

hasta:

\[ x_f=\SI{11}{\meter}, \]

sin cambiar de sentido.

Determine:

\begin{enumerate}[label=\alph*)]
    \item el desplazamiento;
    \item la distancia recorrida.
\end{enumerate}

\textit{Pista: si no existe cambio de sentido, la distancia coincide con la magnitud del desplazamiento.}
\end{ejercicio}

\begin{respuesta}
\[
\Delta x=+\SI{9}{\meter},
\qquad
d=\SI{9}{\meter}.
\]
\end{respuesta}

\begin{explicacion}
El desplazamiento es:

\[ \Delta x=x_f-x_i. \]

Por tanto:

\[
\Delta x
=
\SI{11}{\meter}
-
\SI{2}{\meter}
=
+\SI{9}{\meter}.
\]

Así:

\[ \boxed{\Delta x=+\SI{9}{\meter}}. \]

Como la persona se mueve directamente desde \(x=2\) hasta \(x=11\) sin cambiar de sentido, recorre exactamente \(\SI{9}{\meter}\):

\[ \boxed{d=\SI{9}{\meter}}. \]

En este caso:

\[ d=|\Delta x|. \]

Esto ocurre porque no hubo retrocesos ni cambios de sentido.
\end{explicacion}


\begin{ejercicio}
Una estudiante parte del origen de una pista recta, camina hasta:

\[ x=+\SI{8}{\meter}, \]

y después regresa hasta:

\[ x=+\SI{3}{\meter}. \]

Determine:

\begin{enumerate}[label=\alph*)]
    \item la posición final;
    \item el desplazamiento total;
    \item la distancia total recorrida.
\end{enumerate}

\textit{Pista: para la distancia debe sumar la longitud de cada tramo recorrido.}
\end{ejercicio}

\begin{respuesta}
\[
x_f=+\SI{3}{\meter},
\qquad
\Delta x=+\SI{3}{\meter},
\qquad
d=\SI{13}{\meter}.
\]
\end{respuesta}

\begin{explicacion}
La posición inicial es:

\[ x_i=0. \]

La posición final es:

\[ \boxed{x_f=+\SI{3}{\meter}}. \]

El desplazamiento depende únicamente de las posiciones inicial y final:

\[
\Delta x=x_f-x_i
=
\SI{3}{\meter}-0.
\]

Por tanto:

\[ \boxed{\Delta x=+\SI{3}{\meter}}. \]

Para determinar la distancia debemos analizar el recorrido.

Primer tramo:

\[ d_1=\SI{8}{\meter}. \]

Segundo tramo, desde \(x=8\) hasta \(x=3\):

\[ d_2=\SI{5}{\meter}. \]

Entonces:

\[
d=d_1+d_2
=
\SI{8}{\meter}
+
\SI{5}{\meter}.
\]

Así:

\[ \boxed{d=\SI{13}{\meter}}. \]

La diferencia entre ambos resultados muestra que distancia y desplazamiento no representan la misma magnitud física.
\end{explicacion}


\begin{ejercicio}
Un experimento comienza en el instante:

\[ t_i=\SI{3.5}{\second} \]

y termina en:

\[ t_f=\SI{11.2}{\second}. \]

Determine el intervalo de tiempo transcurrido.

\textit{Pista: \(\Delta t=t_f-t_i\).}
\end{ejercicio}

\begin{respuesta}
\[ \boxed{\Delta t=\SI{7.7}{\second}}. \]
\end{respuesta}

\begin{explicacion}
El intervalo de tiempo se define como:

\[ \Delta t=t_f-t_i. \]

Sustituyendo:

\[
\Delta t
=
\SI{11.2}{\second}
-
\SI{3.5}{\second}.
\]

Entonces:

\[ \boxed{\Delta t=\SI{7.7}{\second}}. \]

El instante \(t_i\) indica cuándo comienza el intervalo y \(t_f\) cuándo termina. El intervalo \(\Delta t\) representa cuánto tiempo transcurrió entre ambos acontecimientos.
\end{explicacion}


\begin{ejercicio}
Clasifique las siguientes trayectorias como \textbf{rectilíneas} o \textbf{curvilíneas}:

\begin{enumerate}[label=\alph*)]
    \item un ascensor que sube por un conducto vertical recto;
    \item un automóvil que recorre una curva de una carretera;
    \item una piedra que se mueve por una guía horizontal recta;
    \item un punto del borde de una rueda mientras esta gira alrededor de un eje fijo.
\end{enumerate}

\textit{Pista: la trayectoria es el conjunto de posiciones sucesivas ocupadas por el objeto.}
\end{ejercicio}

\begin{respuesta}
\begin{enumerate}[label=\alph*)]
    \item Rectilínea.
    \item Curvilínea.
    \item Rectilínea.
    \item Curvilínea.
\end{enumerate}
\end{respuesta}

\begin{explicacion}
Una trayectoria es rectilínea cuando todas las posiciones sucesivas del objeto se encuentran sobre una línea recta.

Por tanto:

\begin{enumerate}[label=\alph*)]
    \item El ascensor se mueve sobre una línea vertical: trayectoria rectilínea.

    \item El automóvil cambia continuamente la dirección geométrica de su recorrido: trayectoria curvilínea.

    \item La piedra permanece sobre una guía recta: trayectoria rectilínea.

    \item El punto del borde de la rueda describe una curva alrededor del eje: trayectoria curvilínea.
\end{enumerate}

La clasificación se refiere a la forma geométrica de la trayectoria, no a que el objeto tenga rapidez constante o variable.
\end{explicacion}


\begin{ejercicio}
Indique en cuáles de las siguientes situaciones resulta razonable modelar el objeto como una \textbf{partícula}:

\begin{enumerate}[label=\alph*)]
    \item estudiar el recorrido de un automóvil entre dos ciudades separadas por \(\SI{300}{\kilo\meter}\);
    \item estudiar si el automóvil cabe dentro de un espacio de estacionamiento;
    \item estudiar el movimiento de la Tierra alrededor del Sol;
    \item estudiar la rotación de la Tierra alrededor de su propio eje.
\end{enumerate}

Justifique brevemente.

\textit{Pista: pregúntese si el tamaño, la forma o la orientación del cuerpo son importantes para el problema.}
\end{ejercicio}

\begin{respuesta}
El modelo de partícula es adecuado en a) y c), pero no resulta suficiente en b) ni d).
\end{respuesta}

\begin{explicacion}
Modelar un cuerpo como partícula significa representar su posición mediante un solo punto y despreciar sus dimensiones y orientación.

\begin{enumerate}[label=\alph*)]
    \item Para estudiar un viaje de cientos de kilómetros, las dimensiones del automóvil son pequeñas respecto de la distancia recorrida. El modelo de partícula es razonable.

    \item Para saber si cabe en un estacionamiento, la longitud y el ancho del automóvil son esenciales. No puede tratarse simplemente como un punto.

    \item Para estudiar la posición de la Tierra a lo largo de su órbita, su tamaño puede ser pequeño comparado con las dimensiones del sistema solar. El modelo de partícula puede ser útil.

    \item Para estudiar la rotación terrestre es indispensable considerar su extensión y orientación. El modelo de partícula no contiene esa información.
\end{enumerate}

El modelo adecuado depende de la pregunta física que se desea responder.
\end{explicacion}


\nivelintermedio


\begin{ejercicio}
Una persona se mueve sobre una línea recta. Parte de:

\[ x_i=-\SI{5}{\meter}. \]

Luego realiza el siguiente recorrido:

\[
-\SI{5}{\meter}
\longrightarrow
+\SI{7}{\meter}
\longrightarrow
-\SI{2}{\meter}
\longrightarrow
+\SI{4}{\meter}.
\]

Determine:

\begin{enumerate}[label=\alph*)]
    \item la posición final;
    \item el desplazamiento total;
    \item la distancia total recorrida.
\end{enumerate}

\textit{Pista: calcule por separado la longitud de cada tramo antes de sumarlos.}
\end{ejercicio}

\begin{respuesta}
\[
x_f=+\SI{4}{\meter},
\qquad
\Delta x=+\SI{9}{\meter},
\qquad
d=\SI{27}{\meter}.
\]
\end{respuesta}

\begin{explicacion}
La posición inicial es:

\[ x_i=-\SI{5}{\meter}. \]

La posición final es:

\[ x_f=+\SI{4}{\meter}. \]

Por tanto, el desplazamiento es:

\[
\Delta x
=
x_f-x_i
=
\SI{4}{\meter}
-
(-\SI{5}{\meter}).
\]

Entonces:

\[ \boxed{\Delta x=+\SI{9}{\meter}}. \]

Para la distancia analizamos cada tramo.

Primer tramo:

\[
d_1
=
|7-(-5)|
=
\SI{12}{\meter}.
\]

Segundo tramo:

\[
d_2
=
|-2-7|
=
\SI{9}{\meter}.
\]

Tercer tramo:

\[
d_3
=
|4-(-2)|
=
\SI{6}{\meter}.
\]

Por tanto:

\[
d=d_1+d_2+d_3
=
12+9+6.
\]

Así:

\[ \boxed{d=\SI{27}{\meter}}. \]

Aunque la persona recorrió \(\SI{27}{\meter}\), su cambio neto de posición fue solamente de \(\SI{9}{\meter}\) hacia el sentido positivo.
\end{explicacion}


\begin{ejercicio}
En un sistema de referencia, una motocicleta se encuentra en:

\[ x=\SI{18}{\meter}. \]

Se decide cambiar el origen y colocarlo en un punto que, en el sistema anterior, tenía coordenada:

\[ x_O=\SI{5}{\meter}. \]

El sentido positivo no cambia.

¿Cuál es la posición de la motocicleta en el nuevo sistema?

\textit{Pista: la nueva coordenada mide la separación algebraica respecto del nuevo origen.}
\end{ejercicio}

\begin{respuesta}
\[ \boxed{x'=\SI{13}{\meter}}. \]
\end{respuesta}

\begin{explicacion}
En el sistema original:

\[ x=\SI{18}{\meter}. \]

El nuevo origen se encuentra en:

\[ x_O=\SI{5}{\meter}. \]

Como el sentido positivo permanece igual, la nueva coordenada es:

\[ x'=x-x_O. \]

Entonces:

\[
x'
=
\SI{18}{\meter}
-
\SI{5}{\meter}.
\]

Por tanto:

\[ \boxed{x'=\SI{13}{\meter}}. \]

El objeto no se movió físicamente. Lo que cambió fue la forma de asignarle una coordenada debido al cambio de origen.
\end{explicacion}


\begin{ejercicio}
En cierto sistema de referencia, una partícula se mueve desde:

\[ x_i=-\SI{3}{\meter} \]

hasta:

\[ x_f=+\SI{8}{\meter}. \]

\begin{enumerate}[label=\alph*)]
    \item Determine el desplazamiento usando el eje original.
    \item Ahora mantenga el mismo origen, pero invierta el sentido positivo del eje. ¿Cuáles serán las nuevas coordenadas \(x_i'\) y \(x_f'\)?
    \item Determine \(\Delta x'\).
    \item Explique por qué los signos de ambos desplazamientos son diferentes aunque se trate del mismo movimiento físico.
\end{enumerate}

\textit{Pista: al invertir el eje, las coordenadas cambian de signo.}
\end{ejercicio}

\begin{respuesta}
\[
\Delta x=+\SI{11}{\meter},
\qquad
x_i'=+\SI{3}{\meter},
\qquad
x_f'=-\SI{8}{\meter},
\qquad
\Delta x'=-\SI{11}{\meter}.
\]
\end{respuesta}

\begin{explicacion}
En el sistema original:

\[
\Delta x=x_f-x_i.
\]

Entonces:

\[
\Delta x
=
\SI{8}{\meter}
-
(-\SI{3}{\meter})
=
+\SI{11}{\meter}.
\]

Al invertir el sentido positivo sin cambiar el origen:

\[
x_i'=+\SI{3}{\meter},
\]

y:

\[
x_f'=-\SI{8}{\meter}.
\]

Por tanto:

\[
\Delta x'
=
x_f'-x_i'
=
-\SI{8}{\meter}
-
\SI{3}{\meter}.
\]

Así:

\[ \boxed{\Delta x'=-\SI{11}{\meter}}. \]

El movimiento físico es exactamente el mismo. Lo que cambia es la componente algebraica utilizada para describirlo.

En el primer sistema, el desplazamiento tiene el sentido positivo; en el segundo, ese mismo sentido físico ha sido definido como negativo.

Esto muestra que los signos dependen del sistema de referencia elegido.
\end{explicacion}


\begin{ejercicio}
La posición de un móvil se registra solamente en ciertos instantes:

\begin{center}
\begin{tabular}{c|cccc}
\toprule
\(t\) (\si{\second}) & 0 & 1 & 2 & 3 \\
\midrule
\(x\) (\si{\meter}) & 2 & 5 & 5 & 1 \\
\bottomrule
\end{tabular}
\end{center}

Determine:

\begin{enumerate}[label=\alph*)]
    \item la posición inicial;
    \item la posición final;
    \item el desplazamiento entre \(t=0\) y \(t=\SI{3}{\second}\);
    \item la distancia mínima que necesariamente tuvo que recorrer;
    \item si puede asegurarse que permaneció en reposo entre \(t=\SI{1}{\second}\) y \(t=\SI{2}{\second}\).
\end{enumerate}

\textit{Pista: conocer la misma posición en dos instantes no revela necesariamente qué ocurrió entre ellos.}
\end{ejercicio}

\begin{respuesta}
\begin{enumerate}[label=\alph*)]
    \item \(x_i=\SI{2}{\meter}\).
    \item \(x_f=\SI{1}{\meter}\).
    \item \(\Delta x=-\SI{1}{\meter}\).
    \item \(d_{\min}=\SI{7}{\meter}\).
    \item No puede asegurarse.
\end{enumerate}
\end{respuesta}

\begin{explicacion}
La posición inicial es:

\[ x_i=\SI{2}{\meter}, \]

y la final:

\[ x_f=\SI{1}{\meter}. \]

El desplazamiento total es:

\[
\Delta x
=
x_f-x_i
=
1-2.
\]

Por tanto:

\[ \boxed{\Delta x=-\SI{1}{\meter}}. \]

Para determinar la menor distancia compatible con los datos, suponemos que el móvil recorrió directamente la separación entre cada par de posiciones registradas.

De \(x=2\) a \(x=5\):

\[ d_1=\SI{3}{\meter}. \]

Entre \(t=1\) y \(t=2\), las posiciones registradas son ambas \(x=5\). La distancia mínima posible es:

\[ d_2=0. \]

De \(x=5\) a \(x=1\):

\[ d_3=\SI{4}{\meter}. \]

Así:

\[ \boxed{d_{\min}=\SI{7}{\meter}}. \]

Sin embargo, que:

\[ x(\SI{1}{\second})=x(\SI{2}{\second}) \]

no demuestra que el móvil haya permanecido en reposo. Podría haberse alejado de \(x=5\) y regresar antes de la segunda medición.

Por ello, la información discreta de posición no permite reconstruir necesariamente toda la trayectoria.
\end{explicacion}


\begin{ejercicio}
Una atleta corre una vuelta completa por una pista cerrada de \(\SI{400}{\meter}\) y termina exactamente en el punto donde comenzó.

Determine:

\begin{enumerate}[label=\alph*)]
    \item la distancia recorrida;
    \item la magnitud del desplazamiento;
    \item si la trayectoria es rectilínea o curvilínea.
\end{enumerate}

\textit{Pista: el desplazamiento depende de las posiciones inicial y final, no de la longitud del recorrido.}
\end{ejercicio}

\begin{respuesta}
\[
d=\SI{400}{\meter},
\qquad
|\Delta\vect{r}|=0.
\]

La trayectoria es curvilínea y cerrada.
\end{respuesta}

\begin{explicacion}
La atleta completa toda la longitud de la pista:

\[ \boxed{d=\SI{400}{\meter}}. \]

Sin embargo, termina exactamente en la misma posición en la que comenzó.

Por definición:

\[ \Delta\vect{r}=\vect{r}_f-\vect{r}_i. \]

Como:

\[ \vect{r}_f=\vect{r}_i, \]

resulta:

\[ \boxed{\Delta\vect{r}=\vect{0}}. \]

Su magnitud es:

\[ \boxed{|\Delta\vect{r}|=0}. \]

La trayectoria no es una línea recta, sino una curva cerrada.

Este ejemplo muestra claramente que un objeto puede recorrer una distancia grande y, al mismo tiempo, tener desplazamiento nulo.
\end{explicacion}


\begin{ejercicio}
En una carretera recta se elige un árbol como origen y se toma el este como sentido positivo.

Una casa \(A\) se encuentra en:

\[ x_A=-\SI{35}{\meter}, \]

y una tienda \(B\) en:

\[ x_B=+\SI{20}{\meter}. \]

Una persona camina directamente desde \(A\) hasta \(B\).

Determine:

\begin{enumerate}[label=\alph*)]
    \item su desplazamiento;
    \item la distancia recorrida;
    \item el sentido del desplazamiento.
\end{enumerate}

\textit{Pista: ambas posiciones están a lados opuestos del origen.}
\end{ejercicio}

\begin{respuesta}
\[
\Delta x=+\SI{55}{\meter},
\qquad
d=\SI{55}{\meter}.
\]

El desplazamiento apunta hacia el este.
\end{respuesta}

\begin{explicacion}
La posición inicial es:

\[ x_i=-\SI{35}{\meter}, \]

y la final:

\[ x_f=+\SI{20}{\meter}. \]

Entonces:

\[
\Delta x
=
x_f-x_i
=
20-(-35).
\]

Por tanto:

\[ \boxed{\Delta x=+\SI{55}{\meter}}. \]

Como la persona camina directamente entre ambos puntos y no cambia de sentido:

\[ \boxed{d=\SI{55}{\meter}}. \]

El signo positivo del desplazamiento indica que su sentido es el definido como positivo, es decir, hacia el este.
\end{explicacion}


\begin{ejercicio}
En un sistema de referencia se registra:

\[
x_i=\SI{12}{\meter},
\qquad
x_f=-\SI{4}{\meter}.
\]

Después se cambia el origen al punto que anteriormente tenía coordenada:

\[ x_O=\SI{7}{\meter}, \]

sin modificar el sentido positivo.

\begin{enumerate}[label=\alph*)]
    \item Determine el desplazamiento en el sistema original.
    \item Determine las posiciones inicial y final en el nuevo sistema.
    \item Calcule nuevamente el desplazamiento.
    \item ¿Cambió el desplazamiento al trasladar el origen?
\end{enumerate}

\textit{Pista: restar la misma cantidad a ambas posiciones no cambia su diferencia.}
\end{ejercicio}

\begin{respuesta}
\[
\Delta x=-\SI{16}{\meter},
\]

\[
x_i'=\SI{5}{\meter},
\qquad
x_f'=-\SI{11}{\meter},
\]

\[
\Delta x'=-\SI{16}{\meter}.
\]

El desplazamiento no cambia.
\end{respuesta}

\begin{explicacion}
En el sistema original:

\[
\Delta x=x_f-x_i.
\]

Entonces:

\[
\Delta x
=
-\SI{4}{\meter}
-
\SI{12}{\meter}.
\]

Por tanto:

\[ \boxed{\Delta x=-\SI{16}{\meter}}. \]

Al trasladar el origen hasta \(x_O=\SI{7}{\meter}\):

\[
x_i'=x_i-x_O
=
12-7
=
\SI{5}{\meter},
\]

y:

\[
x_f'=x_f-x_O
=
-4-7
=
-\SI{11}{\meter}.
\]

El nuevo desplazamiento es:

\[
\Delta x'
=
x_f'-x_i'
=
-11-5.
\]

Así:

\[ \boxed{\Delta x'=-\SI{16}{\meter}}. \]

El desplazamiento permanece igual porque trasladar el origen modifica ambas coordenadas en la misma cantidad.

La posición depende de dónde colocamos el origen; el cambio de posición entre dos puntos no depende de esa traslación.
\end{explicacion}


\begin{ejercicio}
Dos personas observan a una pasajera sentada dentro de un autobús que avanza por una carretera.

La primera persona está sentada junto a ella dentro del autobús. La segunda se encuentra de pie en la acera.

\begin{enumerate}[label=\alph*)]
    \item ¿La pasajera está en reposo o en movimiento respecto de la primera persona?
    \item ¿Está en reposo o en movimiento respecto de la segunda?
    \item ¿Existe contradicción entre ambas respuestas?
\end{enumerate}

\textit{Pista: el movimiento debe describirse siempre respecto de un sistema de referencia.}
\end{ejercicio}

\begin{respuesta}
\begin{enumerate}[label=\alph*)]
    \item Está en reposo respecto del observador dentro del autobús.
    \item Está en movimiento respecto del observador en la acera.
    \item No existe contradicción.
\end{enumerate}
\end{respuesta}

\begin{explicacion}
Respecto del autobús, la posición de la pasajera sentada no cambia. Por tanto, para el observador que viaja junto a ella se encuentra en reposo.

Respecto de la carretera, la posición de la pasajera cambia continuamente a medida que el autobús avanza. Por tanto, para el observador situado en la acera está en movimiento.

Ambas afirmaciones son correctas porque se refieren a sistemas de referencia diferentes.

\begin{interpretacion}
Reposo y movimiento no son propiedades absolutas de un objeto.

Siempre deben especificarse respecto de qué sistema de referencia se está describiendo el movimiento.
\end{interpretacion}
\end{explicacion}


\nivelavanzado


\begin{ejercicio}
Diseñe dos movimientos rectilíneos que comiencen en:

\[ x_i=0 \]

y terminen en:

\[ x_f=+\SI{5}{\meter}, \]

de modo que:

\begin{enumerate}[label=\alph*)]
    \item en el primer movimiento la distancia recorrida sea \(\SI{5}{\meter}\);
    \item en el segundo la distancia recorrida sea \(\SI{13}{\meter}\).
\end{enumerate}

Determine el desplazamiento en ambos casos y explique qué demuestra el resultado acerca de la diferencia entre distancia y desplazamiento.
\end{ejercicio}

\begin{respuesta}
Un ejemplo posible es:

\[
0\longrightarrow5
\]

para el primer caso, y:

\[
0\longrightarrow-4\longrightarrow5
\]

para el segundo.

En ambos:

\[ \Delta x=+\SI{5}{\meter}. \]

Las distancias son, respectivamente:

\[ \SI{5}{\meter}\quad\text{y}\quad\SI{13}{\meter}. \]
\end{respuesta}

\begin{explicacion}
En ambos movimientos:

\[
x_i=0,
\qquad
x_f=\SI{5}{\meter}.
\]

Por tanto, el desplazamiento es necesariamente:

\[
\Delta x=x_f-x_i
=
\SI{5}{\meter}.
\]

Así:

\[ \boxed{\Delta x=+\SI{5}{\meter}}. \]

Para el primer recorrido podemos movernos directamente:

\[ 0\longrightarrow5. \]

La distancia es:

\[ d_1=\SI{5}{\meter}. \]

Para obtener una distancia mayor sin cambiar las posiciones inicial y final podemos introducir un cambio de sentido:

\[ 0\longrightarrow-4\longrightarrow5. \]

La primera parte mide:

\[ \SI{4}{\meter}. \]

La segunda:

\[ \SI{9}{\meter}. \]

Entonces:

\[
d_2
=
4+9
=
\SI{13}{\meter}.
\]

Por tanto, dos movimientos pueden tener exactamente el mismo desplazamiento y distancias recorridas diferentes.

El desplazamiento contiene información sobre el cambio neto de posición; la distancia contiene información sobre la longitud total del recorrido.
\end{explicacion}


\begin{ejercicio}
Demuestre, para un movimiento rectilíneo, que la distancia recorrida \(d\) nunca puede ser menor que la magnitud del desplazamiento:

\[ d\geq|\Delta x|. \]

Indique además en qué situación se cumple la igualdad.
\end{ejercicio}

\begin{respuesta}
Siempre:

\[ \boxed{d\geq|\Delta x|}. \]

La igualdad se cumple cuando el objeto se mueve directamente entre las posiciones inicial y final sin invertir el sentido.
\end{respuesta}

\begin{explicacion}
Supongamos que el movimiento se divide en varios desplazamientos parciales:

\[
\Delta x_1,\Delta x_2,\ldots,\Delta x_n.
\]

El desplazamiento total es:

\[
\Delta x
=
\Delta x_1+\Delta x_2+\cdots+\Delta x_n.
\]

La distancia recorrida es la suma de las magnitudes de esos cambios:

\[
d
=
|\Delta x_1|
+
|\Delta x_2|
+\cdots+
|\Delta x_n|.
\]

Por la desigualdad triangular:

\[
|\Delta x_1+\Delta x_2+\cdots+\Delta x_n|
\leq
|\Delta x_1|
+
|\Delta x_2|
+\cdots+
|\Delta x_n|.
\]

Por tanto:

\[ \boxed{|\Delta x|\leq d}. \]

La igualdad ocurre cuando todos los desplazamientos parciales tienen el mismo sentido, es decir, cuando el objeto no retrocede.

Si existen cambios de sentido, parte de los desplazamientos algebraicos se cancela al calcular \(\Delta x\), pero todas las longitudes recorridas se suman positivamente en \(d\). Por ello:

\[ d>|\Delta x|. \]
\end{explicacion}


\begin{ejercicio}
De un móvil sólo se conocen los siguientes datos:

\[
x(0)=\SI{0}{\meter},
\qquad
x(\SI{10}{\second})=\SI{0}{\meter}.
\]

Analice cuáles de las siguientes afirmaciones pueden asegurarse:

\begin{enumerate}[label=\alph*)]
    \item El desplazamiento durante los \(\SI{10}{\second}\) fue cero.
    \item La distancia recorrida fue cero.
    \item El móvil permaneció en reposo durante todo el intervalo.
    \item El móvil pudo haberse desplazado y luego regresar al origen.
    \item La distancia recorrida pudo ser mayor que cero.
\end{enumerate}

Justifique cada respuesta.
\end{ejercicio}

\begin{respuesta}
Son necesariamente verdaderas a), d) y e) en el sentido de posibilidad para d) y e). No pueden asegurarse b) ni c).
\end{respuesta}

\begin{explicacion}
La posición inicial y la final son iguales:

\[
x_i=x_f=0.
\]

Por tanto:

\[
\Delta x=x_f-x_i=0.
\]

Así, a) es necesariamente verdadera.

Sin embargo, este resultado no permite determinar la trayectoria intermedia.

Una posibilidad es que el objeto haya permanecido siempre en:

\[ x=0. \]

En ese caso:

\[ d=0. \]

Pero también podría haber realizado, por ejemplo:

\[
0
\longrightarrow
+\SI{20}{\meter}
\longrightarrow
0.
\]

Entonces:

\[
d
=
\SI{40}{\meter},
\]

mientras que:

\[ \Delta x=0. \]

Por tanto, no podemos asegurar que la distancia sea cero ni que el móvil haya permanecido en reposo.

Sí podemos afirmar que es posible que se haya movido y regresado, y que es perfectamente posible que la distancia recorrida sea mayor que cero.

Este ejercicio muestra que conocer únicamente dos posiciones no permite reconstruir todo el movimiento entre ellas.
\end{explicacion}


\begin{ejercicio}
Un estudiante afirma:

\begin{quote}
``Si la posición de un objeto es negativa, entonces necesariamente se está moviendo en el sentido negativo del eje''.
\end{quote}

Analice críticamente la afirmación y construya al menos dos situaciones diferentes que demuestren si es correcta o incorrecta.
\end{ejercicio}

\begin{respuesta}
La afirmación es incorrecta.

Un objeto con \(x<0\) puede moverse hacia el sentido positivo, hacia el negativo o incluso estar momentáneamente en reposo.
\end{respuesta}

\begin{explicacion}
La posición y el movimiento describen aspectos diferentes.

La condición:

\[ x<0 \]

significa únicamente que el objeto se encuentra en el lado negativo del origen.

Por ejemplo, considere un objeto en:

\[ x=-\SI{6}{\meter}. \]

Primer caso: puede estar desplazándose hacia la derecha, es decir, hacia posiciones mayores. En ese instante se encuentra en una posición negativa, pero su movimiento tiene sentido positivo.

Segundo caso: el mismo objeto puede desplazarse hacia la izquierda. En ese caso tanto su posición como el sentido de su movimiento son negativos.

Tercer caso: también podría encontrarse momentáneamente en reposo en:

\[ x=-\SI{6}{\meter}. \]

Por tanto, el signo de la posición no permite determinar el sentido del movimiento.

\begin{interpretacion}
La coordenada responde a la pregunta:

\begin{quote}
¿Dónde está el objeto?
\end{quote}

El estado de movimiento requiere información adicional acerca de cómo cambia esa posición.
\end{interpretacion}
\end{explicacion}


\begin{ejercicio}
Se desea estudiar a una gimnasta durante un salto con giro.

Analice si el modelo de partícula es adecuado para cada uno de los siguientes objetivos:

\begin{enumerate}[label=\alph*)]
    \item describir aproximadamente cómo se traslada la gimnasta de un extremo al otro del escenario;
    \item determinar cuántas vueltas completas realiza durante el salto;
    \item estudiar si sus pies y sus manos ocupan determinadas posiciones durante una figura;
    \item registrar únicamente la trayectoria de un punto representativo de su cuerpo.
\end{enumerate}

Explique qué información se pierde cuando se utiliza el modelo de partícula.
\end{ejercicio}

\begin{respuesta}
El modelo de partícula puede ser adecuado para a) y d), pero no para b) ni c).

Al utilizarlo se pierde información sobre tamaño, forma, orientación y rotación del cuerpo.
\end{respuesta}

\begin{explicacion}
Una partícula posee posición, pero no tiene extensión espacial ni orientación.

Para estudiar únicamente la traslación global de la gimnasta entre dos regiones del escenario, puede bastar con representar su cuerpo mediante un punto.

Del mismo modo, si sólo interesa la trayectoria de un punto representativo, el modelo puede resultar útil.

En cambio, para determinar cuántas vueltas realiza es necesario describir la orientación del cuerpo y su rotación. Esa información desaparece en el modelo de partícula.

También es inadecuado si necesitamos conocer simultáneamente dónde se encuentran manos y pies, porque eso exige considerar dimensiones y forma.

Este ejemplo muestra que un modelo no es correcto o incorrecto en términos absolutos. Su utilidad depende de la pregunta física que queremos responder.
\end{explicacion}


\begin{ejercicio}
Dos observadores utilizan sistemas de coordenadas diferentes para describir el mismo movimiento rectilíneo.

El observador \(A\) utiliza un eje \(x\). El observador \(B\) utiliza otro eje \(x'\) cuyo origen se encuentra \(\SI{20}{\meter}\) a la derecha del origen de \(A\), pero ambos ejes tienen el mismo sentido positivo.

Un automóvil se mueve desde:

\[ x_i=\SI{8}{\meter} \]

hasta:

\[ x_f=\SI{35}{\meter} \]

según el observador \(A\).

Determine:

\begin{enumerate}[label=\alph*)]
    \item las posiciones inicial y final medidas por \(B\);
    \item el desplazamiento medido por \(A\);
    \item el desplazamiento medido por \(B\);
    \item qué magnitudes cambian y cuáles permanecen iguales al trasladar el origen.
\end{enumerate}
\end{ejercicio}

\begin{respuesta}
\[
x_i'=-\SI{12}{\meter},
\qquad
x_f'=+\SI{15}{\meter},
\]

y:

\[
\Delta x=\Delta x'=+\SI{27}{\meter}.
\]

Las posiciones cambian; el desplazamiento no.
\end{respuesta}

\begin{explicacion}
El nuevo origen se encuentra en:

\[ x_O=\SI{20}{\meter}. \]

Como los sentidos positivos son iguales:

\[ x'=x-x_O. \]

Para la posición inicial:

\[
x_i'
=
\SI{8}{\meter}
-
\SI{20}{\meter}
=
-\SI{12}{\meter}.
\]

Para la posición final:

\[
x_f'
=
\SI{35}{\meter}
-
\SI{20}{\meter}
=
+\SI{15}{\meter}.
\]

Por tanto:

\[
\boxed{x_i'=-\SI{12}{\meter}},
\qquad
\boxed{x_f'=+\SI{15}{\meter}}.
\]

El observador \(A\) mide:

\[
\Delta x
=
x_f-x_i
=
35-8
=
+\SI{27}{\meter}.
\]

El observador \(B\) obtiene:

\[
\Delta x'
=
x_f'-x_i'
=
15-(-12)
=
+\SI{27}{\meter}.
\]

Así:

\[ \boxed{\Delta x=\Delta x'=+\SI{27}{\meter}}. \]

Las coordenadas individuales dependen de la ubicación elegida para el origen. En cambio, la diferencia entre las posiciones inicial y final permanece igual cuando únicamente trasladamos el origen y mantenemos la orientación del eje.

Este resultado distingue claramente entre \textbf{posición} y \textbf{desplazamiento}.
\end{explicacion}